\documentclass[12pt]{article}
\usepackage{amsmath, amssymb}
\usepackage{geometry}
\geometry{margin=1in}
\usepackage{setspace}
\usepackage{enumitem}

\begin{document}
\onehalfspacing

\begin{center}
\textbf{CSE400 – Fundamentals of Probability in Computing} \\
\textbf{Lecture 9 Scribe} \\
\textbf{Uniform, Exponential, Laplace and Gamma Random Variables} \\
Instructor: Dhaval Patel, PhD \\
Date: February 2, 2026
\end{center}

\section{Introduction and Lecture Scope}
This lecture introduces types of continuous random variables and discusses their mathematical formulation, properties, and applications. The lecture specifically covers:

\begin{itemize}
\item Uniform Random Variable
\item Exponential Random Variable
\item Laplace Random Variable
\item Gamma Random Variable
\item Graphical and example-based understanding
\item Problem solving and Gaussian density estimation discussion
\end{itemize}

The material establishes the mathematical definitions of distributions using Probability Density Function (PDF) and Cumulative Distribution Function (CDF). L9\_S2\_A

\section{Types of Continuous Random Variables}
Continuous random variables are defined using probability density functions instead of probability mass functions. The lecture focuses on commonly used continuous distributions used in computing and modelling scenarios. L9\_S2\_A

\section{Uniform Random Variable}

\subsection{Definition}
A random variable $X$ is called a Uniform Random Variable over the interval $[a,b]$ if the probability density function is constant over this interval and zero elsewhere.

\subsection{Probability Density Function (PDF)}
The PDF of a uniform random variable is given by:
\[
f_X(x) =
\begin{cases}
\frac{1}{b-a}, & a \leq x < b \\
0, & \text{elsewhere}
\end{cases}
\]

\textbf{Explanation}

The probability is evenly distributed across the interval $[a,b]$.

Since probability density must integrate to 1:
\[
\int_a^b f_X(x)\,dx = 1
\]

Because the density is constant over the interval:
\[
f_X(x) = \frac{1}{b-a}
\]

\subsection{Cumulative Distribution Function (CDF)}
The CDF is obtained by integrating the PDF.
\[
F_X(x) =
\begin{cases}
0, & x < a \\
\frac{x-a}{b-a}, & a \leq x < b \\
1, & x \geq b
\end{cases}
\]

\subsection{Step-by-Step Derivation of CDF}

\textbf{Case 1:} $x < a$

Since the distribution starts at $a$,
\[
F_X(x) = 0
\]

\textbf{Case 2:} $a \leq x < b$

\[
F_X(x) = \int_a^x \frac{1}{b-a} \, dt
\]

Evaluating integral:
\[
F_X(x) = \frac{x-a}{b-a}
\]

\textbf{Case 3:} $x \geq b$

Entire probability mass is accumulated:
\[
F_X(x) = 1
\]

\subsection{Observations from Lecture Discussion}
The lecture emphasizes:

\begin{itemize}
\item Uniform distribution is often treated as a reference distribution.
\item It is used in modelling and simulation.
\item It forms a base distribution in many algorithms.
\end{itemize}

\section{Exponential Random Variable}

\subsection{Definition}
An exponential random variable models time-related processes such as waiting time and decay-based behaviour.

\subsection{Probability Density Function}
For parameter $b > 0$, the PDF is defined as:
\[
f_X(x) = \frac{1}{b} \exp\left(-\frac{x}{b}\right) u(x)
\]

Where:

\begin{itemize}
\item $u(x)$ is the unit step function
\item $u(x) = 0$ for $x < 0$
\item $u(x) = 1$ for $x \geq 0$
\end{itemize}

\textbf{Interpretation of Unit Step Function}

The unit step function restricts the domain of the exponential random variable such that:
\[
X \geq 0
\]

This means the distribution exists only for non-negative values.

\subsection{Cumulative Distribution Function}
The CDF is given by:
\[
F_X(x) = \left(1 - \exp\left(-\frac{x}{b}\right)\right) u(x)
\]

\subsection{Step-by-Step Derivation of CDF}

The CDF is computed using integration of PDF.

\textbf{Step 1: Start with definition}
\[
F_X(x) = \int_{-\infty}^{x} f_X(t)\,dt
\]

\textbf{Step 2: Apply support restriction}

Since exponential RV exists only for $t \geq 0$:
\[
F_X(x) =
\begin{cases}
0, & x < 0 \\
\int_0^x \frac{1}{b} e^{-t/b} \, dt, & x \geq 0
\end{cases}
\]

\textbf{Step 3: Solve integral}

\[
\int_0^x \frac{1}{b} e^{-t/b} \, dt
\]

Let:
\[
I = \int e^{-t/b} \, dt
\]

\[
I = -b e^{-t/b}
\]

Applying limits:
\[
F_X(x) = \left[-e^{-t/b}\right]_0^x
\]

\[
F_X(x) = 1 - e^{-x/b}
\]

\textbf{Step 4: Include unit step function}

\[
F_X(x) = (1 - e^{-x/b})u(x)
\]

\subsection{Graphical Behaviour}

\textbf{PDF Characteristics}

\begin{itemize}
\item Maximum value occurs at $x = 0$
\item PDF decreases exponentially as $x$ increases
\item Distribution is skewed to the right
\end{itemize}

\textbf{CDF Characteristics}

\begin{itemize}
\item Starts at zero
\item Monotonically increases
\item Approaches 1 asymptotically
\end{itemize}

\section{Example – Exponential Random Variable}

\textbf{Problem Statement}

Let $X$ be an exponential random variable with PDF:
\[
f_X(x) = e^{-x}u(x)
\]

\textbf{Step 1: Identify Distribution Support}

From unit step function:
\[
X \geq 0
\]

\textbf{Step 2: Interpret Expression}

The presence of $u(x)$ ensures:

\begin{itemize}
\item The PDF is zero for negative values.
\item The exponential decay starts from $x = 0$.
\end{itemize}

\textbf{Step 3: Observed Domain Representation}

The lecture explains the domain as:
\[
X \in (-\infty, \infty)
\]

However, due to $u(x)$, the effective domain becomes:
\[
X \in [0, \infty)
\]

\section{Laplace Random Variable}
The lecture lists the Laplace random variable as part of the covered distributions. The discussion includes:

\begin{itemize}
\item Example explanation
\item Distribution usage in modelling
\item Graph-based understanding
\end{itemize}

\section{Gamma Random Variable}
The lecture introduces Gamma random variables and includes:

\begin{itemize}
\item Graph and special case analysis
\item Example explanation
\item Homework problem related to Gamma distribution
\end{itemize}

\section{Applications and Contextual Discussion}
The lecture links continuous random variables with:

\begin{itemize}
\item Randomised algorithms
\item Modelling processes
\item Coding and inference mechanisms
\item Data modelling and estimation
\item Gaussian density estimation (in-class activity)
\end{itemize}

\section{Key Learning Observations}

\begin{itemize}
\item Continuous distributions are defined using PDF and CDF.
\item Uniform distribution models equal likelihood across intervals.
\item Exponential distribution models decay and waiting-time processes.
\item Unit step functions are used to restrict distribution domains.
\item Graphical interpretation helps understand behaviour of distributions.
\item These distributions are widely used in computing, modelling, and algorithmic design.
\end{itemize}

\section{Summary}
This lecture develops foundational understanding of major continuous random variable distributions. The lecture systematically derives:

\begin{itemize}
\item Mathematical definitions
\item PDF formulations
\item CDF derivations
\item Graphical interpretations
\item Application contexts
\end{itemize}

The material forms a basis for probabilistic modelling and statistical analysis in computing systems. L9\_S2\_A

\end{document}
